\section{Penance, Confession, and the Law of the Season}\label{sec:penance-law}

On the first weekday morning of Advent---Monday of the First Week---the Roman Breviary opens Isaiah at Matins with a summons to purification: \textit{Lavamini, mundi estote}, ``Wash yourselves, be clean.'' Before two days of the season have passed, the Church places before her children the Lord's promise in Isaiah~1:16--18:

\begin{quote}
Wash yourselves, be clean, take away the evil of your devices from my eyes, cease to do perversely, Learn to do well: seek judgment, relieve the oppressed, judge for the fatherless, defend the widow. And then come, and accuse me, saith the Lord: if your sins be as scarlet, they shall be made as white as snow: and if they be red as crimson, they shall be white as wool.
\end{quote}

Throughout Advent, the Church pairs sacramental confession with the outward discipline of fasting, self-denial, and mercy.

\subsection{Confession: Making Straight His Paths}

As Christmas draws near, the Gospel of Ember Saturday and the Fourth Sunday of Advent brings the Church down to the Jordan to hear St.~John the Baptist ``preaching the baptism of penance for the remission of sins'' (Luke~3:3) and echoing Isaiah's cry, ``Prepare ye the way of the Lord, make straight his paths'' (Luke~3:4; Section~\ref{sec:week-4}). John refuses easy assurances and demands conversion visible in conduct: ``Bring forth therefore fruits worthy of penance'' (Luke~3:8)---bidding the crowds share coat and food, tax collectors exact nothing beyond what is appointed, and soldiers put away violence and extortion (Luke~3:10--14).

Plan to make one thorough, unhurried Advent confession. The natural week in the Church's calendar is the Ember week---between Gaudete Sunday and the Great O~Antiphons, when the liturgy itself turns to fasting and penitential prayer---or earlier in Advent if your conscience asks for it. Going before the final week spares you the haste of Christmas Eve lines and lets you enter the Great O~Antiphons and the Vigil in peace.

In his \textit{Introduction to a Devout Life} (Part~II, Chapter~19), St.~Francis de Sales presented confession not as a dreaded tribunal, but as a sovereign healing:

\begin{quote}
Our Saviour has left the holy sacrament of penance and confession to his Church, that in it we might cleanse ourselves from all our iniquities, as often as we should be defiled by them. Never suffer your heart then, O~Philothea! to remain long infected with sin, since you have so easy a remedy at hand.
\end{quote}

Recommending frequent confession---a spiritual counsel to Philothea rather than a precept of church law---he explained why a Christian should come even without mortal sin on the conscience: ``by confession you not only receive absolution from the venial sins you confess, but likewise strength to avoid them, light to discern them well, and grace to repair all the damage you may have sustained by them \ldots{} in this one act of confession you shall exercise more virtues than in any other whatsoever.'' He warned against vague formulas---``Beware of unmeaning self-accusations, made out of a mere routine''---and urged instead ``a hearty sorrow for the sins you confess, however small they are; as also a stedfast resolution to correct them in future.''

A serene way to prepare both nightly and before Advent confession is the five-point General Examen from the \textit{Spiritual Exercises} of St.~Ignatius of Loyola. Ignatius begins with gratitude rather than anxiety, because sins appear most clearly beside God's mercy:

\begin{enumerate}
\item \textbf{Give thanks.} ``The first Point is to give thanks to God our Lord for the benefits received.''
\item \textbf{Ask for light.} ``The second, to ask grace to know our sins and cast them out.''
\item \textbf{Review your life.} ``The third, to ask account of our soul from the hour that we rose up to the present Examen, hour by hour, or period by period: and first as to thoughts, and then as to words, and then as to acts, in the same order as was mentioned in the Particular Examen.'' Before confession, widen this review across the months since your last confession against the Commandments and the duties of your state of life.
\item \textbf{Ask pardon.} ``The fourth, to ask pardon of God our Lord for the faults.''
\item \textbf{Purpose amendment.} ``The fifth, to purpose amendment with His grace,'' closing with the ``Our Father.''
\end{enumerate}

In the confessional, speak plainly: state every mortal sin in kind and number, name concrete venial faults you mean to amend, receive absolution in faith, and perform your penance promptly.

\subsection{Fasting, Abstinence, Precept, and Closed Times: Then and Now}

Church discipline changes across centuries and territories while the divine call to penance abides.

\paragraph{Earlier historical discipline.}
As seen in Section~\ref{sec:history}, sixth-century Gaul kept St.~Martin's Lent from 11~November to Christmas---three fasts a week under Bishop Perpetuus of Tours and the First Council of M\^acon (581--583), and daily December fasting for monks under the Second Council of Tours (567)---while the Council of Trent (Session~XXIV, 11~November~1563) enforced for the Latin Church the ban on solemn marriages from Advent Sunday until Epiphany.

Before universal codification in 1917, stricter regional fasts persisted. In England, Dom Laurence Shepherd noted in 1870 that Advent Wednesdays and Fridays remained days of fast and abstinence. In the United States, the \textit{Manual of Prayers} ordered by the Third Plenary Council of Baltimore (1884; published 1889) listed as fasting days under precept ``The Fridays of Advent'' (noting provincial exemptions such as New Orleans), the Advent Ember Days, and the Vigil of Christmas---anticipating any Sunday fast to Saturday---while keeping Trent's closed time for solemnizing marriage ``From Advent Sunday until the Epiphany.''

\paragraph{Universal Latin law in 1962.}
In 1962, universal Latin discipline rested on the 1917 \textit{Codex Iuris Canonici} (promulgated 27~May~1917, effective 19~May~1918, superseded 27~November~1983) and a universal decree of 1959. Every Friday obliged to abstinence alone (canon~1252, paragraph~1); outside the Ember week and Christmas Eve, the 1917 Code imposed no universal Advent fast and no Advent Friday fast. Fast and abstinence together obliged on the Advent Ember Days and the Vigil of Christmas (canon~1252, paragraph~2); both ceased on Sundays and holy days of precept outside Lent, and vigils were no longer anticipated to Saturday when falling on Sunday (canon~1252, paragraph~4). On 3~December~1959, Pope St.~John~XXIII, through the Sacred Congregation of the Council, granted all the faithful the universal faculty to anticipate the Christmas Vigil fast and abstinence from 24~December to 23~December. For marriage, canon~1108 allowed matrimony at any time of year (paragraph~1), forbade only the solemn nuptial blessing from Advent~I through Christmas Day inclusive (paragraph~2, shortening Trent's Epiphany boundary), and authorized the local Ordinary to permit it for a just cause without excessive pomp (paragraph~3).

\paragraph{Current universal Latin law.}
On 17~February~1966, Pope St.~Paul~VI promulgated \textit{Paenitemini} (Norms~II, IV, and~VI), restricting universal fast and abstinence to Ash Wednesday and Good Friday---ending statutory Ember-Day and Christmas-Vigil fasts---while keeping Friday abstinence and empowering episcopal conferences to substitute works of charity and piety.

Every Latin Catholic today, including those attending the 1962 Roman Rite, lives under the 1983 \textit{Codex Iuris Canonici} (promulgated 25~January~1983 by \textit{Sacrae disciplinae leges}; effective First Sunday of Advent, 27~November~1983). Sunday is the primordial holy day of obligation alongside ten universal holy days, including the Immaculate Conception (8~December) and the Nativity of the Lord (25~December); the faithful must participate in Mass and abstain from works hindering worship, joy, or rest, subject to conference suppressions or Sunday transfers approved by the Apostolic See (canons~1246--1247). All Catholics must do penance by divine law (canon~1249), and universal statutory penitential days are Fridays and Lent alone (canon~1250): Advent is not a statutory penitential season. Friday abstinence obliges unless a solemnity falls on Friday; fast and abstinence oblige on Ash Wednesday and Good Friday alone (canon~1251). Abstinence obliges from the completed fourteenth year, and fasting from majority until the beginning of the sixtieth year, while pastors and parents form younger children in penance (canon~1252). Episcopal conferences may specify observance further or substitute charity and piety (canon~1253). No general canon in the 1983 Code forbids the nuptial blessing in Advent.

\paragraph{Current United States particular law.}
Under \textit{Paenitemini}, the National Conference of Catholic Bishops issued its \textit{Pastoral Statement on Penance and Abstinence} on 18~November~1966, effective First Sunday of Advent, 27~November~1966. Retained under the 1983 Code by the USCCB Complementary Norm to canons~1252--1253 (approved September~1983; promulgated 21~October~1983) and amended 29~February~1984, United States law sets abstinence from age fourteen (canon~1252) and fasting ``from the completion of the eighteenth year to the beginning of the sixtieth.'' Advent carries no statutory fast or abstinence: ``Vigils and Ember Days, as most now know, no longer oblige to fast and abstinence'' (paragraph~17), and Advent is kept by ``meditation on the lessons taught by the liturgy and of increased participation in the liturgical rites'' (paragraph~7). Friday remains ``a special day of penitential observance'' (paragraph~22), yet outside Lent abstinence from meat no longer obliges ``under pain of sin,'' though the bishops ``give first place to abstinence from flesh meat'' among voluntary Friday penances (paragraphs~24--25). For holy-day obligation rules in any given year and country---including when 8~December falls on Saturday, Sunday, or Monday---consult your diocese and bishops' conference.

\begin{longtable}{>{\raggedright\arraybackslash}p{0.17\linewidth}>{\raggedright\arraybackslash}p{0.24\linewidth}>{\raggedright\arraybackslash}p{0.25\linewidth}>{\raggedright\arraybackslash}p{0.25\linewidth}}
\toprule
\textbf{Observance} & \textbf{Earlier History (Gaul, Trent, U.S.~1889)} & \textbf{1962 Universal Law (1917~CIC \& 1959 Decree)} & \textbf{Present Law (1983~CIC \& U.S.~Norms)} \\
\midrule
\endhead
\textbf{Advent \& Fridays} & Gaul (M\^acon c.~9, 581/583): Mon/Wed/Fri fast 11~Nov--25~Dec. U.S.~(1889): Advent Fridays fast (with diocesan exemptions). & Friday abstinence only (1917~CIC c.~1252 \S 1); no universal Advent fast outside Ember Days and Vigil. & Universal: Advent not statutory penitential season (c.~1250); Friday abstinence unless solemnity (c.~1251). U.S.: Voluntary Friday penance outside Lent; meat abstinence first (NCCB 1966 \S\S 22--25). \\
\textbf{Ember Days} & U.S.~(1889): Fast and abstinence Wed, Fri, Sat after Advent~III. & Fast and abstinence Wed, Fri, Sat (1917~CIC c.~1252 \S 2). & No statutory fast or abstinence (\textit{Paenitemini} 1966; 1983~CIC cc.~1250--1251; NCCB 1966 \S 17). \\
\textbf{Christmas Vigil (24~Dec)} & U.S.~(1889): Fast and abstinence; kept Saturday if Sunday. & Fast and abstinence (c.~1252 \S 2); not anticipated if Sunday (c.~1252 \S 4); faculty to anticipate to 23~Dec (3~Dec~1959, AAS~51:918). & No statutory fast or abstinence (\textit{Paenitemini} 1966; 1983~CIC cc.~1250--1251; NCCB 1966 \S 17). \\
\textbf{Holy days (8~Dec, 25~Dec)} & Regional calendars varied; 25~Dec lifted Friday abstinence (1889, p.~18). & Fast and abstinence ceased on Sundays and holy days outside Lent (c.~1252 \S 4). & 8~Dec and 25~Dec among ten universal holy days (cc.~1246--1247); consult diocese and conference for local rules. \\
\textbf{Closed times} & Solemn marriages forbidden Advent~I--Epiphany (Trent Sess.~XXIV c.~10, 1563; 1889, p.~19). & Solemn nuptial blessing forbidden Advent~I--25~Dec inclusive; Ordinary could permit for just cause (c.~1108 \S\S 1--3). & No general closed-time canon in the 1983 Code. \\
\bottomrule
\end{longtable}

Full canon citations and jurisdictional bounds appear in Appendix~\ref{app:about}.

\paragraph{Precept and generous voluntary penance.}
That church law no longer commands an Advent fast clarifies the line between sin and generosity. A Catholic in the United States who eats meat on an Advent Friday, or does not fast on Ember Wednesday or Christmas Eve, breaks no law of the Church. Yet in paragraph~17 of the 1966 statement, the bishops encouraged the devout to prepare freely for great feasts ``by a day of particular self-denial, penitential prayer and fasting.'' St.~Francis de Sales offered the same balanced counsel (\textit{Introduction to the Devout Life}, Part~III, Chapter~23): where health and duties permit, ``observe some days beyond what are ordered by the Church, for besides the ordinary effect of fasting in raising the mind, subduing the flesh, confirming goodness, and obtaining a heavenly reward,'' voluntary restraint keeps appetite subject to the spirit.

Keep that restraint proportionate to health and household: abstain from meat on Advent Fridays, take a simpler meal on the Ember Days or Christmas Eve, quiet evening noise, and give what is spared to the poor. Voluntary penance in Advent serves wakefulness, as Pope St.~Leo the Great urged in his nineteenth sermon on the December fast (section~1):

\begin{quote}
For the advent of which it behoves every man to prepare himself, lest it find him given over to gluttony, or entangled in cares of this life.
\end{quote}
